\chapter{1981年全国卷（理）}

\section{解答题}

\begin{problem}
设 \(A\) 表示有理数的集合，\(B\) 表示无理数的集合，即设 \(A = \{\text{有理数}\}\)，\(B = \{\text{无理数}\}\)，试写出：

\begin{enumerate}
\item \(A \cup B\)；
\item \(A \cap B\).
\end{enumerate}
\end{problem}

\begin{solution}
\begin{enumerate}
\item 若 \(x\in A\cup B\)，则 \(x\) 是有理数或无理数，因而 \(x\in\R\)；反过来，每个实数或者是有理数，或者是无理数，所以 \(\R\subseteq A\cup B\). 因此 \(A\cup B=\R\).
\item 若 \(x\in A\cap B\)，则 \(x\) 同时是有理数和无理数，这是不可能的，所以不存在这样的 \(x\)，从而 \(A\cap B=\varnothing\).
\end{enumerate}
\end{solution}

\begin{problem}
在 A，B，C，D 四位候选人中，

\begin{enumerate}
\item 如果选举正，副班长各一人，共有几种选法？写出所有可能的选举结果；
\item 如果选举班委三人，共有几种选法？写出所有可能的选举结果.
\end{enumerate}
\end{problem}

\begin{solution}
\begin{enumerate}
\item 先选正班长，有 \(4\) 种选法；再从剩下的 \(3\) 人中选副班长，有 \(3\) 种选法，所以共有 \(4\times3=12\) 种. 按正班长、
副班长的顺序写出全部结果为 AB，AC，AD，BC，BD，CD，BA，CA，DA，CB，DB，DC.
\item 三名班委只区分人选而不区分职务，因此从四人中任取三人，选法数为 \(\mathrm C_4^3=4\). 全部结果为 ABC，ABD，ACD，BCD.
\end{enumerate}
\end{solution}

\begin{problem}
下表所列各小题中，指出 \(A\) 是 \(B\) 的充分条件，还是必要条件，还是充要条件，或者都不是.

\begin{center}
\renewcommand{\arraystretch}{1.25}
\setlength{\tabcolsep}{8pt}
\begin{tabular}{|c|c|c|}
\hline
\(A\) & \(B\) & \(A\) 是 \(B\) 的什么条件 \\
\hline
四边形 \(ABCD\) 为平行四边形 & 四边形 \(ABCD\) 为矩形 & \\
\hline
\(a=3\) & \(|a|=3\) & \\
\hline
\(\theta=150^\circ\) & \(\sin\theta=1/2\) & \\
\hline
点 \((a,b)\) 在圆 \(x^{2}+y^{2}=R^{2}\) 上 & \(a^{2}+b^{2}=R^{2}\) & \\
\hline
\end{tabular}
\end{center}
\end{problem}

\begin{solution}
\begin{enumerate}
\item 矩形一定是平行四边形，所以由条件 \(B\) 可以推出条件 \(A\)，即 \(A\) 是 \(B\) 的必要条件. 平行四边形不一定是矩形，故不是充分条件.
\item 由 \(a=3\) 可得 \(|a|=3\)，所以 \(A\) 是 \(B\) 的充分条件；但 \(a=-3\) 时 \(|a|=3\) 而 \(a\ne3\)，故不是必要条件.
\item 由 \(\theta=150^\circ\) 得 \(\sin\theta=\sin30^\circ=1/2\)，所以 \(A\) 是 \(B\) 的充分条件；例如 \(\theta=30^\circ\) 时也有 \(\sin\theta=1/2\)，但不满足 \(\theta=150^\circ\)，故 \(A\) 不是必要条件.
\item 点 \((a,b)\) 在圆 \(x^2+y^2=R^2\) 上的定义就是其坐标满足 \(a^2+b^2=R^2\)，所以两个条件可以相互推出，\(A\) 是 \(B\) 的充要条件.
\end{enumerate}
\end{solution}

\begin{problem}
写出余弦定理（只写一个公式即可），并加以证明.
\end{problem}

\begin{solution}
设 \(\triangle ABC\) 的三个内角 \(A\)，\(B\)，\(C\) 的对边分别为 \(a\)，\(b\)，\(c\). 以点 \(A\) 为原点，射线 \(AB\) 为 \(x\) 轴正向建立直角坐标系，则
\(B=(c,0)\)，\(C=(b\cos A,b\sin A)\). 由两点间距离公式，
\[
a^2=|BC|^2=(c-b\cos A)^2+(0-b\sin A)^2=b^2+c^2-2bc\cos A
\]
这就是余弦定理. 坐标表示对锐角、直角和钝角都成立，故证明完毕.
\end{solution}

\begin{problem}
解不等式（\(x\) 为未知数）：\( \left|\begin{matrix}x-a & b & -c \\ a & x-b & c \\ -a & b & x-c\end{matrix}\right| > 0 \).
\end{problem}

\begin{solution}
记左端的行列式为 \(\Delta\). 将第二行加到第一行，再将第三行加到第二行，行列式的值不变，于是
\[
\Delta=\begin{vmatrix}x&x&0\\0&x&x\\-a&b&x-c\end{vmatrix}
=x^2\begin{vmatrix}1&1&0\\0&1&1\\-a&b&x-c\end{vmatrix}
=x^2\bigl((x-c)-b-a\bigr)=x^2(x-a-b-c)
\]
因此原不等式等价于 \(x^2(x-a-b-c)>0\). 由于 \(x^2\geqslant0\)，要使乘积严格为正，必须且只需 \(x^2>0\) 且 \(x-a-b-c>0\)，所以解为 \(x\ne0\) 且 \(x>a+b+c\).
\end{solution}

\begin{problem}
用数学归纳法证明等式： \(\cos \frac{x}{2} \cdot \cos \frac{x}{2^2} \cdot \cos \frac{x}{2^3} \cdot \cdots \cdot \cos \frac{x}{2^n} = \frac{\sin x}{2^n \sin \frac{x}{2^n}}\). 对一切自然数 \(n\) 都成立.
\end{problem}

\begin{solution}
对任意固定的实数 \(x\)，先证明等式
\[
2^n\sin\frac{x}{2^n}\prod_{j=1}^{n}\cos\frac{x}{2^j}=\sin x
\]
对一切自然数 \(n\) 成立.

当 \(n=1\) 时，二倍角公式给出
\[
2\sin\frac{x}{2}\cos\frac{x}{2}=\sin x
\]
所以结论对 \(n=1\) 成立.

假设当 \(n=k\) 时结论成立，即
\[
2^k\sin\frac{x}{2^k}\prod_{j=1}^{k}\cos\frac{x}{2^j}=\sin x
\]
由二倍角公式 \(2\sin\dfrac{x}{2^{k+1}}\cos\dfrac{x}{2^{k+1}}=\sin\dfrac{x}{2^k}\)，有
\[
2^{k+1}\sin\frac{x}{2^{k+1}}\prod_{j=1}^{k+1}\cos\frac{x}{2^j}
=2^k\sin\frac{x}{2^k}\prod_{j=1}^{k}\cos\frac{x}{2^j}=\sin x
\]
所以结论对 \(n=k+1\) 也成立. 由数学归纳法，乘积等式对一切自然数 \(n\) 成立.

当 \(\sin\dfrac{x}{2^n}\ne0\) 时，两边除以 \(2^n\sin\dfrac{x}{2^n}\)，得到题中的等式. 当该正弦值为零时，题中等式的右端没有定义；因此原等式应在其右端有定义的范围内理解.
\end{solution}

\begin{problem}
设 1980 年底我国人口以 \(10\text{ 亿}\) 计算.

\begin{enumerate}
\item 如果我国人口每年比上年平均递增 \(2\%\)，那么到 2000 年底将达到多少？
\item 要使 2000 年底我国人口不超过 \(12\text{ 亿}\)，那么每年比上年平均递增率最高是多少？
\end{enumerate}

\begin{center}
\fbox{%
\begin{minipage}{9.3cm}
\centering
下列对数值可供选用：

{
\setlength{\tabcolsep}{4pt}
\scriptsize
\begin{tabular}{@{}ccc@{}}
\(\lg 1.0087 = 0.00377\) & \(\lg 1.0092 = 0.00396\) & \(\lg 1.0096 = 0.00417\) \\
\(\lg 1.0200 = 0.00860\) & \(\lg 1.2000 = 0.07918\) & \(\lg 1.3098 = 0.11720\) \\
\(\lg 1.4568 = 0.16340\) & \(\lg 1.4859 = 0.17200\) & \(\lg 1.5157 = 0.18060\)
\end{tabular}
}
\end{minipage}%
}
\end{center}
\end{problem}

\begin{solution}
\begin{enumerate}
\item 从 1980 年底到 2000 年底共有 \(20\) 个年度增长期. 设人口数为 \(x\) 亿，则
\(x=10(1.02)^{20}\). 两边取常用对数，得
\[
\lg x=\lg10+20\lg1.0200=1+20\times0.00860=1.17200
\]
由所给数值 \(\lg1.4859=0.17200\)，可得 \(x=10^{1.17200}=10\times1.4859=14.859\). 所以到 2000 年底人口将达到 \(14.859\text{ 亿}\).
\item 设每年平均递增率为小数 \(r\). 人口不超过 \(12\) 亿的条件为
\(10(1+r)^{20}\leqslant12\)，即 \((1+r)^{20}\leqslant1.2\). 因为对数函数在正数上单调递增，
\(20\lg(1+r)\leqslant\lg1.2000=0.07918\)，\(\lg(1+r)\leqslant0.003959\).
因此精确的最高递增率应为
\[
r_{\max}=\sqrt[20]{1.2}-1
\]
数值计算得 \(r_{\max}\approx0.0091578\)，即约为 \(0.9158\%\). 若按题目给出的对数表核对，
\(\lg1.0087=0.00377<0.003959<0.00396=\lg1.0092\)，所以
\(1.0087<1+r_{\max}<1.0092\)，按百分数保留两位小数，最高递增率为 \(0.92\%\).
\end{enumerate}
\end{solution}

\begin{problem}
在 \(120^{\circ}\) 的二面角 \(P - a - Q\) 的两个面 \(P\) 和 \(Q\) 内，分别有点 \(A\) 和点 \(B\)，已知点 \(A\) 和点 \(B\) 到棱 \(a\) 的距离分别为 \(2\) 和 \(4\)，且线段 \(AB = 10\).

\begin{enumerate}
\item 求直线 \(AB\) 和棱 \(a\) 所成的角；
\item 求直线 \(AB\) 和平面 \(Q\) 所成的角.
\end{enumerate}
\end{problem}

\begin{solution}
\begin{enumerate}
\item 在平面 \(P\) 内作 \(AD\perp a\)，垂足为 \(D\)；在平面 \(Q\) 内作 \(BE\perp a\)，垂足为 \(E\). 在平面 \(Q\) 内过 \(D\) 作 \(DC\perp a\)，并过 \(B\) 作 \(BC\parallel a\)，两线交于 \(C\). 于是四边形 \(BCDE\) 为矩形，从而 \(CD=BE=4\). 又 \(\angle ADC\) 是二面角的平面角，所以 \(\angle ADC=120^\circ\). 在 \(\triangle ACD\) 中，余弦定理给出
\[
AC^2=AD^2+CD^2-2AD\cdot CD\cos120^\circ=2^2+4^2-2\cdot2\cdot4\left(-\frac12\right)=28
\]
即 \(AC=2\sqrt7\). 由于 \(AD\perp a\) 且 \(CD\perp a\)，平面 \(ACD\perp a\)；又 \(BC\parallel a\)，所以 \(BC\perp AC\). 因此 \(\triangle ABC\) 在 \(C\) 处为直角三角形，且 \(BC\parallel a\)，故直线 \(AB\) 与棱 \(a\) 所成的角等于 \(\angle ABC\). 于是
\[
\sin\angle ABC=\frac{AC}{AB}=\frac{2\sqrt7}{10}=\frac{\sqrt7}{5}
\]
所以所求角为 \(\arcsin\dfrac{\sqrt7}{5}\).
\item 平面 \(ACD\) 垂直于棱 \(a\)，而 \(a\subset Q\)，所以平面 \(ACD\perp Q\). 在平面 \(ACD\) 内过 \(A\) 作 \(AF\perp CD\). 因为 \(\angle ADC=120^\circ\)，垂足 \(F\) 位于 \(CD\) 在 \(D\) 的反向延长线上，故 \(\angle ADF=60^\circ\)，从而
\[
AF=AD\sin60^\circ=2\cdot\frac{\sqrt3}{2}=\sqrt3
\]
由平面垂直的性质，\(AF\perp Q\). 又 \(BF\subset Q\)，所以 \(AF\perp BF\)，\(\triangle ABF\) 在 \(F\) 处为直角三角形. 直线 \(AB\) 与平面 \(Q\) 所成的角为 \(\angle ABF\)，故
\[
\sin\angle ABF=\frac{AF}{AB}=\frac{\sqrt3}{10}
\]
所以所求角为 \(\arcsin\dfrac{\sqrt3}{10}\).
\end{enumerate}
\end{solution}

\begin{problem}
给定双曲线 \(x^{2}-\frac{y^{2}}{2}=1\).

\begin{enumerate}
\item 过点 \(A(2,1)\) 的直线 \(L\) 与所给的双曲线交于两点 \(P_{1}\) 及 \(P_{2}\)，求线段 \(P_{1}P_{2}\) 的中点 \(P\) 的轨迹方程；
\item 过点 \(B(1,1)\) 能否作直线 \(m\)，使 \(m\) 与所给双曲线交于两点 \(Q_{1}\) 及 \(Q_{2}\)，且点 \(B\) 是线段 \(Q_{1}Q_{2}\) 的中点？ 这样的直线 \(m\) 如果存在，求出它的方程；如果不存在，说明理由.
\end{enumerate}
\end{problem}

\begin{solution}
\begin{enumerate}
\item 设弦 \(P_1P_2\) 的中点为 \(P(u,v)\)，并设弦的方向向量为 \((p,q)\). 因为 \(P_1\ne P_2\)，可写成
\(P_1=(u+tp,v+tq)\)，\(P_2=(u-tp,v-tq)\)，\(t\ne0\).
把两点分别代入双曲线方程 \(x^2-y^2/2=1\)，相减得
\[
2t\left(2up-vq\right)=0
\]
所以 \(2up-vq=0\). 当 \(P\ne A\) 时，弦所在直线经过 \(A(2,1)\)，可取方向向量 \((p,q)=(u-2,v-1)\)，从而
\[
2u(u-2)-v(v-1)=0
\]
即
\[
v^2-v+2u(2-u)=0
\]
将 \((u,v)\) 改记为轨迹变量 \((x,y)\)，得到必要条件
\[
y^2-y+2x(2-x)=0
\]

下面证明这个方程上的每一点都能取得. 先取 \(P=A=(2,1)\)，直线 \(y=4x-7\) 与双曲线联立得
\[
14x^2-56x+51=0
\]
该方程的判别式为 \(280>0\)，且两根之和为 \(4\)，所以两个交点横坐标的平均值为 \(2\). 由于交点都在直线 \(y=4x-7\) 上，它们纵坐标的平均值为 \(4\times2-7=1\)，故两个交点的中点为 \((2,1)=A\).

再取 \(P=(2,0)\)，直线 \(x=2\) 与双曲线联立得 \(y^2=6\)，交点为 \((2,\sqrt6)\) 和 \((2,-\sqrt6)\)，其中点确为 \(P\).

对于轨迹方程上的其他点 \(P=(u,v)\)，直线 \(AP\) 不是竖直线，可设其方程为
\[
y=k(x-2)+1
\]
若 \(k=\pm\sqrt2\)，将此直线代入轨迹方程，得到
\[
(x-2)\left[(k-4)+(k^2-2)(x-2)\right]=(k-4)(x-2)=0
\]
由于 \(k\ne4\)，这时轨迹方程在线 \(AP\) 上只有点 \(A\)，与 \(P\ne A\) 矛盾. 因此 \(k^2\ne2\).

将直线代入双曲线，得到关于交点横坐标的二次方程
\[
(2-k^2)x^2+(4k^2-2k)x-4k^2+4k-3=0
\]
其二次项系数不为零，且判别式为
\[
(4k^2-2k)^2-4(2-k^2)(-4k^2+4k-3)
=8(3k^2-4k+3)=8\left(3\left(k-\frac23\right)^2+\frac53\right)>0
\]
所以该直线与双曲线交于两个不同的有限点. 设两交点横坐标为 \(x_1\)，\(x_2\)，由韦达定理，交点横坐标的中点为
\[
\frac{x_1+x_2}{2}=-\frac{2k^2-k}{2-k^2}
\]
另一方面，点 \(P=(u,v)\) 在轨迹方程和直线 \(AP\) 上. 因为 \(P\ne A\)，代入轨迹方程后可得
\[
(u-2)\left[(k-4)+(k^2-2)(u-2)\right]=0
\]
故
\[
u=2+\frac{4-k}{k^2-2}=-\frac{2k^2-k}{2-k^2}
\]
于是 \(u=(x_1+x_2)/2\). 又因两个交点都在直线 \(y=k(x-2)+1\) 上，它们纵坐标的中点为
\[
k\left(\frac{x_1+x_2}{2}-2\right)+1=k(u-2)+1=v
\]
因此 \(P\) 确为弦中点. 所求轨迹方程为
\[
y^2-y+2x(2-x)=0
\]
\item 若 \(m\) 为竖直线 \(x=1\)，它与双曲线只有切点 \((1,0)\)，不可能交于两点. 因此设 \(m\) 的方程为 \(y=k(x-1)+1\). 代入双曲线方程并整理，得
\[
(2-k^2)x^2+(2k^2-2k)x-k^2+2k-3=0
\]
要有两个有限交点，必须有 \(2-k^2\ne0\). 设两个交点的横坐标为 \(x_1\)，\(x_2\). 因为 \(B(1,1)\) 是弦中点，所以 \(x_1+x_2=2\). 由韦达定理，
\[
x_1+x_2=-\frac{2k^2-2k}{2-k^2}=2
\]
从而 \(k=2\). 但将 \(k=2\) 代入上面的二次方程，得到
\[
-2x^2+4x-3=0
\]
其判别式为 \(4^2-4(-2)(-3)=-8<0\)，没有实交点. 因此不存在满足条件的直线 \(m\).
\end{enumerate}
\end{solution}

\section{附加题}

\begin{problem}
已知以 \(AB\) 为直径的半圆有一个内接正方形 \(CDEF\)，其边长为 \(1\)（如图）. 设 \(AC = a\)，\(BC = b\)，作数列 \(u_1 = a - b\)，\(u_2 = a^2 - ab + b^2\)，\(u_3 = a^3 - a^2b + ab^2 - b^3\)，\(\cdots\)，\(u_k = a^k - a^{k-1}b + a^{k-2}b^2 - \cdots + (-1)^k b^k\)；求证： \(u_n = u_{n-1} + u_{n-2}\)（\(n \geqslant 3\)）.

\bitmapfigure[max width=4.6cm]{img/1981/national_science/q_additional_fig1.png}
\end{problem}

\begin{solution}
由正方形和半圆的对称性，圆心 \(O\) 是线段 \(FC\) 的中点，所以 \(OF=OC=\dfrac12\). 又 \(CD=1\)，且 \(OD\) 是半圆半径，故在直角三角形 \(OCD\) 中
\[
OA^2=OD^2=OC^2+CD^2=\frac14+1=\frac54
\]
从而 \(AB=2OA=\sqrt5\). 同时 \(AF=BC\)，且 \(AC=AF+FC\)，因此
\(a+b=AB=\sqrt5\)，\(a-b=AC-BC=FC=1\).
于是 \(ab=\dfrac{(a+b)^2-(a-b)^2}{4}=1\). 由 \(a-b=1\) 和 \(ab=1\)，有
\(a^2=a(a-b)+ab=a+1\)，\(b=b(a-b)=ab-b^2=1-b^2\)
即 \(b^2=1-b\). 令 \(c=-b\)，则 \(c^2=c+1\).

对任意正整数 \(k\)，题设数列的通项是等比数列求和公式
\[
u_k=a^k-a^{k-1}b+\cdots+(-1)^kb^k
=\frac{a^{k+1}-(-b)^{k+1}}{a+b}
\]
由于 \(a^2=a+1\) 且 \((-b)^2=-b+1\)，分别有
\(a^{n+1}=a^n+a^{n-1}\)，\((-b)^{n+1}=(-b)^n+(-b)^{n-1}\).
因此当 \(n\geqslant3\) 时
\[
u_{n-1}+u_{n-2}
=\frac{a^n-(-b)^n+a^{n-1}-(-b)^{n-1}}{a+b}
=\frac{a^{n+1}-(-b)^{n+1}}{a+b}
=u_n
\]
故 \(u_n=u_{n-1}+u_{n-2}\)（\(n\geqslant3\)）.
\end{solution}
